\documentclass[aspectratio=169]{beamer}
\makeatletter
\def\input@path{{../}}
\makeatother
\usepackage{theme} % Tu STY severo
\usepackage[table]{xcolor} 
\usepackage{fontawesome5} % Para los iconos de Github, Gitlab y Linux
\usepackage{listings}

% --- FIX PARA YAML Y LOS DOS PUNTOS ---
\lstdefinelanguage{YAML}{
  keywords={true,false,null,y,n},
  keywordstyle=\color{gludOrange}\bfseries,
  basicstyle=\color{gludWhite}\ttfamily\small,
  commentstyle=\color{gludLightGray}\itshape,
  stringstyle=\color{gludOrangeLight},
  moredelim=[l][\color{gludOrange}]{\&},
  moredelim=[l][\color{gludOrangeLight}]{\*},
  morestring=[b]',
  morestring=[b]",
  literate={:}{{\textcolor{gludOrangeLight}{:}}}1 % El fix mágico para los dos puntos
           {-}{{\textcolor{gludOrange}{-}}}1,
}

\lstnewenvironment{yamlcode}{
  \lstset{
    language=YAML,
    breaklines=true,
    breakatwhitespace=true,
    frame=none,
    backgroundcolor=\color{gludDarker},
    xleftmargin=10pt,
    aboveskip=0pt,
    belowskip=0pt
  }
}{}

% Metadatos
\title{Integración Continua (CI)}
\subtitle{Semana 9 - Automatización en la vida real}
\author{Juan Carlos Quintero Rubiano (JkVely)}
\date{Noviembre 2025}

\begin{document}

% 1. Portada
\titleframe[Curso Git 101 - GLUD]

% 2. Índice
\contentsframe

% ==========================================
\section{El concepto de CI}
% ==========================================

\begin{frame}{¿Qué es Integración Continua?}
    \begin{orangebox}[Definición Práctica]
        Es la práctica de subir código a un repositorio compartido (como GitHub o GitLab) \highlight{varias veces al día}. Cada vez que subes algo, un servidor revisa automáticamente si tu código sirve o si acabas de romper el proyecto.
    \end{orangebox}

    \vspace{0.3cm}
    \textbf{¿Qué podemos hacer con CI?}
    \begin{itemize}
        \item Compilar el código (asegurar que no hay errores de sintaxis).
        \item Correr pruebas automáticas.
        \item Revisar el estilo del código (Linting).
        \item Auditar vulnerabilidades de seguridad en librerías.
    \end{itemize}
\end{frame}

% ==========================================
\section{Tipos de Pruebas}
% ==========================================

\begin{frame}{Pruebas Unitarias (Unit Tests)}
    Son la primera línea de defensa de tu código.
    
    \vspace{0.3cm}
    \begin{darkbox}[Características]
        \begin{itemize}
            \item Prueban una sola función o clase de manera aislada.
            \item Son \highlight{extremadamente rápidas} (milisegundos).
            \item No usan bases de datos, ni red, ni archivos externos.
        \end{itemize}
    \end{darkbox}
    
    \vspace{0.3cm}
    \textbf{Ejemplo:} Validar que la función que calcula el IVA en tu carrito de compras devuelva el valor correcto si le pasas $\$100.000$.
\end{frame}

\begin{frame}{Pruebas de Integración}
    Revisan que los bloques de tu sistema se hablen bien entre ellos.
    
    \vspace{0.3cm}
    \begin{alertbox}[Características]
        \begin{itemize}
            \item {\color{black}Prueban el conjunto: Tu código + Base de datos + APIs.}
            \item {\color{black}Son más lentas y costosas de ejecutar.}
            \item {\color{black}Requieren preparar un entorno (ej. levantar contenedores).}
        \end{itemize}
    \end{alertbox}
    
    \vspace{0.3cm}
    \textbf{Ejemplo:} Validar que al registrar un usuario en la página, sus datos efectivamente queden guardados en las tablas de PostgreSQL.
\end{frame}

% ==========================================
% EJEMPLOS PRÁCTICOS DE CÓDIGO
% ==========================================

\begin{frame}[fragile]{Ejemplo: Prueba Unitaria (Java con JUnit)}
    \begin{darkbox}[Contexto: Servicio de Facturaci\'on]
        Probando \textbf{solo} la l\'ogica matem\'atica en memoria, sin tocar la base de datos ni levantar servidores. !Se ejecuta en milisegundos!
    \end{darkbox}

\end{frame}
\begin{frame}[fragile]{Ejemplo: Prueba Unitaria (Java con JUnit)}
    
    \vspace{-0.4cm}
    \begin{codebox}[CalculadoraIvaTest.java]
\begin{lstlisting}[language=Java, basicstyle=\color{gludWhite}\ttfamily\scriptsize, keywordstyle=\color{gludOrange}\bfseries, stringstyle=\color{gludOrangeLight}, commentstyle=\color{gludLightGray}\itshape, backgroundcolor=\color{gludDarker}]
import static org.junit.jupiter.api.Assertions.assertEquals;
import org.junit.jupiter.api.Test;
class CalculadoraIvaTest {
    @Test
    void calcularIva_DebeRetornarEl19Porciento() {
        // Arrange (Preparar datos)
        CalculadoraIva calc = new CalculadoraIva();
        double montoBase = 100000.0;
        // Act (Ejecutar la funcion aislada)
        double resultado = calc.calcularIva(montoBase);
        // Assert (Verificar que el resultado es el esperado)
        assertEquals(19000.0, resultado);
    }
}
\end{lstlisting}
    \end{codebox}
\end{frame}

\begin{frame}[fragile]{Ejemplo: Prueba de Integraci\'on}
    \begin{alertbox}[Contexto: Dashboard de Usuarios]
        Aqu\'i probamos el flujo completo: El navegador renderiza Astro $\rightarrow$ Astro hace fetch a la API $\rightarrow$ La API consulta PostgreSQL.
    \end{alertbox}
\end{frame}
\begin{frame}[fragile]{Ejemplo: Prueba de Integraci\'on}
    \vspace{-0.4cm}
    \begin{codebox}[usuarios.spec.ts (TypeScript)]
\begin{lstlisting}[language=Java, basicstyle=\color{gludWhite}\ttfamily\scriptsize, keywordstyle=\color{gludOrange}\bfseries, stringstyle=\color{gludOrangeLight}, commentstyle=\color{gludLightGray}\itshape, backgroundcolor=\color{gludDarker}]
import { test, expect } from '@playwright/test';

test("Debe cargar la lista de usuarios de la BD", async ({ page }) => {
  // 1. Navegar a la ruta de Astro que hace el SSR (Server Side Rendering)
  await page.goto('http://localhost:4321/usuarios');

  // 2. Esperar que la tabla del DOM se renderice
  const tabla = page.locator('table.usuarios-list');
  await expect(tabla).toBeVisible();

    // 3. Validar que la base de datos devolvio el usuario de prueba
  const nombreUsuario = page.getByText('Juan Perez');
  await expect(nombreUsuario).toBeVisible(); 
    // Si la BD esta caida, esto falla y el CI detiene el pase a produccion
});
\end{lstlisting}
    \end{codebox}
\end{frame}

% ==========================================
\section{Arquitectura YAML}
% ==========================================

\begin{frame}{YAML: El idioma del DevOps}
    \begin{orangebox}[YAML Ain't Markup Language]
        Es un formato súper ligero y legible. Olvídense del XML; YAML es el estándar absoluto para GitHub Actions, GitLab CI, Kubernetes y Ansible.
    \end{orangebox}
    
    \vspace{0.5cm}
    \textbf{Reglas de oro que te salvarán la vida:}
    \begin{itemize}
        \item \highlight{La indentación es la ley:} Usa espacios, \textbf{NUNCA TABS}.
        \item Las listas llevan un guion (\inlinecode{-}).
        \item Los diccionarios (pares) van separados por dos puntos (\inlinecode{:}).
        \item Si indentas mal un espacio, el pipeline explota.
    \end{itemize}
\end{frame}

\begin{frame}{Anatomía de un Pipeline}
    Todo pipeline, sin importar la herramienta, tiene estas partes:
    
    \vspace{0.3cm}
    \begin{itemize}
        \item \step{1}{Triggers:} ¿Qué dispara la acción? (\inlinecode{push}, \inlinecode{pull\_request}, horario).
        \item \step{2}{Runners:} La máquina virtual o contenedor que ejecutará el proceso (Ej: \inlinecode{ubuntu-latest}).
        \item \step{3}{Jobs:} Un grupo de tareas que se ejecutan en un mismo Runner.
        \item \step{4}{Steps:} Los comandos paso a paso (ej. instalar paquetes, compilar, testear).
    \end{itemize}
\end{frame}

% ==========================================
\section{Ecosistema GitHub Actions}
% ==========================================

\begin{frame}[fragile]{\faGithub\ Actions: Triggers y Jobs}
    \begin{codebox}[1. Estructura base (.github/workflows/main.yml)]
\begin{yamlcode}
name: Backend CI
on:
  push:
    branches:
      - main
      - develop

jobs:
  construir-y-probar:
    runs-on: ubuntu-latest
\end{yamlcode}
    \end{codebox}
\end{frame}

\begin{frame}[fragile]{\faGithub\ Actions: Steps y Actions}
    \begin{codebox}[2. Definiendo los pasos]
\begin{yamlcode}
    steps:
      - name: 1. Traer el codigo al servidor
        uses: actions/checkout@v4
    
      - name: 2. Preparar Java 17
        uses: actions/setup-java@v3
        with:
          distribution: 'corretto'
          java-version: '17'

      - name: 3. Ejecutar Maven
        run: mvn clean package
\end{yamlcode}
\end{codebox}
\end{frame}

\begin{frame}{\faGithub\ Las Actions más útiles de la comunidad}
    No reinventen la rueda. Ya hay "plugins" listos para usar:
    
    \vspace{0.3cm}
    \begin{itemize}
        \item \highlight{\inlinecode{actions/checkout@v4}}: Obligatoria. Clona tu repo en la máquina virtual. Sin esto, el runner está vacío.
        \item \highlight{\inlinecode{actions/setup-node}} o \highlight{\inlinecode{setup-java}}: Instalan y configuran el lenguaje de programación y la caché de dependencias.
        \item \highlight{\inlinecode{actions/upload-artifact@v4}}: Saca archivos de la máquina virtual (ej. un \inlinecode{.jar} compilado o un PDF) y los deja descargables en GitHub.
        \item \highlight{\inlinecode{docker/login-action}}: Se autentica en Docker Hub para subir imágenes.
    \end{itemize}
\end{frame}

% ==========================================
\section{Ecosistema GitLab CI}
% ==========================================

\begin{frame}[fragile]{\faGitlab\ GitLab CI: Enfoque en contenedores}
    Todo vive en un solo archivo \inlinecode{.gitlab-ci.yml}. Est\'a dise\~nado alrededor de \faLinux\ Docker.
    
    \begin{codebox}[Estructura base - GitLab CI]
\begin{yamlcode}
image: maven:3.8.5-openjdk-17 # La imagen base de Docker

stages:
  - test
  - build

correr_pruebas:
  stage: test
  script:
    - mvn test
\end{yamlcode}
    \end{codebox}
\end{frame}

% ==========================================
\section{Aplicación Real (GitFlow)}
% ==========================================

\begin{frame}{CI aplicado a GitFlow (El mundo real)}
    ¿Cómo se ve esto en una empresa desarrollando en Spring Boot?
    
    \vspace{0.4cm}
    \begin{itemize}
        \item \step{1}{Rama \inlinecode{feature/*}:} Haces commit. El CI corre el linter y \highlight{Pruebas Unitarias}. Es un proceso rápido que verifica tu lógica básica.
        \item \step{2}{Merge Request a \inlinecode{develop}:} El CI levanta un contenedor de PostgreSQL, arranca tu contexto de Spring y corre \highlight{Pruebas de Integración}. Valida que la base de datos se comunica bien.
        \item \step{3}{Merge a \inlinecode{main}:} Todo está en verde. El CI compila el código, genera el \inlinecode{.jar}, crea una imagen de Docker y la sube al registry.
    \end{itemize}
\end{frame}

% ==========================================
\section{Comparativa Final}
% ==========================================

\begin{frame}{Comparativa directa: \faGithub\ vs \faGitlab}
    \vspace{-0.7cm}
    \begin{comparebox}
        \begin{darkbox}[\faGithub\ GitHub Actions]
            \begin{itemize}
                \item \textbf{Config:} \inlinecode{.github/workflows/*.yml}
                \item \textbf{Reusabilidad:} Marketplace de "Actions".
                \item \textbf{Enfoque:} Orientado a eventos.
                \item \textbf{Curva:} Rápida e intuitiva.
            \end{itemize}
        \end{darkbox}
    \comparebreak
        \begin{darkbox}[\faGitlab\ GitLab CI]
            \begin{itemize}
                \item \textbf{Config:} \inlinecode{.gitlab-ci.yml}
                \item \textbf{Reusabilidad:} Plantillas y scripts puros.
                \item \textbf{Enfoque:} Contenedores Docker nativos.
                \item \textbf{Curva:} Más exigente en bash.
            \end{itemize}
        \end{darkbox}
    \end{comparebox}
    \begin{alertbox}[Conclusión]
    Para sus proyectos, \faGithub\ Actions es más amigable y rápido. \faGitlab\ CI es la herramienta corporativa por excelencia.
    \end{alertbox}
\end{frame}

% ==========================================
\section{Cierre y Manos a la obra}
% ==========================================

\quoteslide{Si tienes que hacer la misma tarea en la terminal más de tres veces a la semana, deberías haberla automatizado en tu CI/CD.}{Ley de la pereza del DevOps}

\thankslide{Gracias\\{\hspace*{1cm}\Large¡A romper producción! (Pero en el CI)}}

\end{document}